\documentclass[10pt,a4paper]{amsart}
\usepackage[margin=19mm]{geometry}
\usepackage{amsmath,amssymb,mathtools}
\usepackage[scr=rsfs,cal=euler]{mathalfa}
\usepackage{xfrac}
\usepackage[unicode,hidelinks,bookmarks=false]{hyperref}
\newcommand{\ie}{i.e.\ }
\newcommand{\cf}{cf.\ }
\newcommand{\resp}{respectively\ }
\newcommand{\Subsection}{\subsection}
\providecommand{\nobreakdash}{\nobreak\hbox{-}}
\setlength{\emergencystretch}{2em}
\multlinegap0pt
\usepackage{bm}
\usepackage{bbm}
\usepackage{dsfont}
\usepackage{xspace}
\usepackage{cancel}
\usepackage{mathtools}
\usepackage{amsthm}
\makeatletter
\@ifundefined{definition}{\newtheorem{definition}{Definition}}{}
\@ifundefined{lemma}{\newtheorem{lemma}{Lemma}}{}
\makeatother
\title[On Time-subordinated Brownian Motion Processes for Financial]{On Time-subordinated Brownian Motion Processes for Financial Markets}
\author{Rohan Shenoy, Peter Kempthorne}
\date{Source: arXiv:2510.14108v2 (CC BY 4.0)}
\begin{document}
\maketitle

\noindent\emph{Source and licence.} On Time-subordinated Brownian Motion Processes for Financial Markets. Version arXiv:2510.14108v2 (CC BY 4.0), distributed under the Creative Commons Attribution 4.0 International licence, as recorded in the arXiv licence metadata for this exact version and shown on the abstract page https://arxiv.org/abs/2510.14108v2. Full rights notice: licenses/arxiv-2510.14108-CC-BY-4.0.txt.\par\smallskip

\noindent\textit{Editorial scope.} Counted unit: the Lemma whose source label is \texttt{ExtendedGaussianCf} in the article (source0.tex lines 444--446, source label \texttt{ExtendedGaussianCf}), delivered together with the article's own complete original proof of it (source0.tex lines 447--470, one \texttt{proof} environment of 24 lines). No mathematics is written, repaired or re-derived: statement and proof are the article's own text line for line. Required context retained uncounted: FourierIntegral (lines 433--439). Every definition, display and label of those lines is retained unchanged. Omitted from this package and not redistributed: 1--432, 440--443, 471--562. Nothing inside a retained excerpt is omitted or altered: every retained body is delivered exactly as the article's lines are written, so no omission receipt of any kind accompanies this package. Citations: the retained excerpts carry no citation command at all, so no bibliography entry of the article is transcribed and no reference is redistributed. Reference adaptation: none was needed and none was made; every reference inside the delivered unit resolves inside it, so no unresolved reference or citation is produced. Printed slips kept literal: no printed word, symbol, hypothesis or number of the counted statement or of its proof was altered. Layout and class: the article's own class is \texttt{article} with packages including authblk, inputenc, graphicx, amsfonts, amsmath, nccmath, amssymb, fancyhdr, amsthm, array, systeme, xcolor, mathtools, afterpage; the wrapper is the lane's pinned \texttt{amsart} editorial wrapper at 10pt on A4, which restates the theorem-like environments the retained text needs and adds the article's own macro definitions that the retained text needs (none), with extra wrapper packages (bm, bbm, dsfont, xspace, cancel, mathtools). The delivered numbering is the unit's own. Assets: no retained span contains a figure environment, a graphics-inclusion command, a PDF-page inclusion, a table or a TikZ picture; no image file is copied into this package, the package declares no asset dependency and no image file is redistributed with it. Licence: version arXiv:2510.14108v2 of this article is distributed under the Creative Commons Attribution 4.0 International licence, as recorded in the arXiv licence metadata for this exact version and shown on the abstract page https://arxiv.org/abs/2510.14108v2. A full rights notice is delivered at licenses/arxiv-2510.14108-CC-BY-4.0.txt. Characters: every retained excerpt is pure ASCII, so the delivered engine, XeLaTeX with its default Latin Modern text font, reports no missing character.
\begin{definition}
    Let $Y$ be a real-valued random variable with density $f_Y(y)$, and $\mathcal{D}\supset\mathbb{R}$ be a region in the complex plane containing the real line such that the Fourier integral
    \begin{equation}\label{FourierIntegral}
        \Psi_Y(z) = \int_{\mathbb{R}}e^{iyz}f_Y(y)dy
    \end{equation}
    is analytic in $\mathcal{D}$. Then $\Psi_Y$ is the analytic continuation of $\psi_Y$ in the domain $\mathcal{D}$.
\end{definition}

\begin{lemma}\label{ExtendedGaussianCf}
    Let $Z$ be a standard Gaussian. Then $\Psi_Z$ is entire in the complex plane $\mathbb{C}$ with $\Psi_Z(z) = e^{-z^2/2}$.
\end{lemma}

\begin{proof}
    Fix $z\in\mathbb{C}$. We begin by completing the square in the Fourier integral (\ref{FourierIntegral}),
    \begin{equation}\label{Fourierintegralstart}
        \int_{\mathbb{R}}e^{ixz}\dfrac{e^{-x^2/2}}{\sqrt{2\pi}}dx = e^{-z^2/2}\int_{\mathbb{R}} \dfrac{1}{\sqrt{2\pi}}e^{-(x-iz)^2/2}dx.
    \end{equation}
    This already has the desired form - once we show that the integral on the RHS is equal to $1$ we may conclude. Let $L>|z|$ and consider the finite integral,
    \begin{equation}
        I_L = \int_{-L}^L \dfrac{1}{\sqrt{2\pi}}e^{-(x-iz)^2/2}dx = \int_{\mathcal{C}_1} \dfrac{1}{\sqrt{2\pi}}e^{-\xi^2/2}d\xi,
    \end{equation}
    where $\mathcal{C}_1$ is the straight line segment $-L-iz\to L-iz$.
    Letting $L\to\infty$, $I_{\infty}$ is the integral we want to evaluate (the integral RHS of (\ref{Fourierintegralstart})). Let $\mathcal{C}$ be the closed loop contour $\mathcal{C} = \mathcal{C}_1\cup \mathcal{C}_2\cup \mathcal{C}_3 \cup \mathcal{C}_4$ encompassing the four line segments (a parallelogram) $$-L-iz\overset{\mathcal{C}_1}{\longrightarrow} L-iz\overset{\mathcal{C}_2}{\longrightarrow} L \overset{\mathcal{C}_3}{\longrightarrow} -L \overset{\mathcal{C}_4}{\longrightarrow} -L-iz.$$ We bound the integrals over $\mathcal{C}_2$ and $\mathcal{C}_4$ using the M-L inequality,
    \begin{align}
        \left|\int_{\mathcal{C}_2}\dfrac{1}{\sqrt{2\pi}}e^{-\xi^2/2}d\xi\right|\leq \dfrac{|z|}{\sqrt{2\pi}}\sup_{\xi\in\mathcal{C}_2}\left|e^{-\xi^2/2}\right|\leq \dfrac{|z|}{\sqrt{2\pi}}e^{-(L-|z|)^2/2},
    \end{align}
    with the same bound for the integral over $\mathcal{C}_4$. Taking the limit as $L\to\infty$, the bound goes to $0$ so that the integrals over $\mathcal{C}_2$ and $\mathcal{C}_4$ vanish as $L\to\infty$.  As $\phi(\xi) = e^{-\xi^2/2}/\sqrt{2\pi}$ is entire, its integral over the closed loop $\mathcal{C}$ is $0$ for all $|L|>z$. This means that in the limit $L\to\infty$,
    \begin{equation}
        \underbrace{\int_{\mathcal{C}_1} \dfrac{1}{\sqrt{2\pi}}e^{-\xi^2/2}d\xi}_{=I_\infty} + 0 + \int_{\mathcal{C}_3}\dfrac{1}{\sqrt{2\pi}}e^{-\xi^2/2}d\xi + 0 = 0.
    \end{equation}
    But for $L\to\infty$ the integral over $\mathcal{C}_3$ is just the Gaussian integral (in the reverse direction),
    \begin{equation}
        \int_{\mathcal{C}_3}\dfrac{1}{\sqrt{2\pi}}e^{-\xi^2/2}d\xi =  \int_{\infty}^{-\infty}\dfrac{1}{\sqrt{2\pi}}e^{-\xi^2/2}d\xi = -1.
    \end{equation}
    So indeed $I_{\infty} = 1$ and we conclude.
\end{proof}

\end{document}
